\documentclass[10pt]{article}
\usepackage[a4paper,margin=0.88in]{geometry}
\usepackage{amsmath,amssymb,amsthm,mathtools,amscd,enumitem,xcolor,hyperref}
\hypersetup{colorlinks=true,linkcolor=blue!55!black,citecolor=blue!55!black,urlcolor=blue!55!black}
\setlength{\emergencystretch}{3em}
\setlength{\parindent}{0pt}
\setlength{\parskip}{0.55em}
\setlist{leftmargin=1.6em,itemsep=0.25em,topsep=0.25em}
\sloppy

\newtheorem{theorem}{Theorem}[section]
\newtheorem{lemma}[theorem]{Lemma}
\newtheorem{proposition}[theorem]{Proposition}
\newtheorem{corollary}[theorem]{Corollary}
\theoremstyle{definition}
\newtheorem{definition}[theorem]{Definition}
\newtheorem{convention}[theorem]{Convention}
\newtheorem{remark}[theorem]{Remark}

\newcommand{\Fq}{\mathbf F_q}
\newcommand{\Fbar}{\overline{\mathbf F}_q}
\newcommand{\Ql}{\mathbf Q_\ell}
\newcommand{\Qlb}{\overline{\mathbf Q}_\ell}
\newcommand{\Dbc}{D^b_c}
\newcommand{\WeilD}{D^b_{\mathrm W,c}}
\newcommand{\WtLe}{\operatorname{WeightLe}}
\newcommand{\WtGe}{\operatorname{WeightGe}}
\newcommand{\IntWtLe}{\operatorname{Wt}^{\mathbf Z}_{\leq}}
\newcommand{\IntWtGe}{\operatorname{Wt}^{\mathbf Z}_{\geq}}
\newcommand{\MixLe}{\operatorname{Mix}^{\leq}}
\newcommand{\MixGe}{\operatorname{Mix}^{\geq}}
\newcommand{\UniLe}{\operatorname{Uni}^{\leq}}
\newcommand{\UniGe}{\operatorname{Uni}^{\geq}}
\newcommand{\Pure}{\operatorname{Pure}}
\newcommand{\RHom}{R\mathcal Hom}
\newcommand{\gr}{\operatorname{gr}}
\newcommand{\Spec}{\operatorname{Spec}}
\newcommand{\AllIntWtLe}{\operatorname{Wt}^{\mathrm{all},\mathbf Z}_{\leq}}
\newcommand{\AllIntWtGe}{\operatorname{Wt}^{\mathrm{all},\mathbf Z}_{\geq}}

\begin{document}
\begin{center}
{\large\bfseries Functoriality of Weights\par}
\medskip
{\normalsize From finite Frobenius modules to the six operations\par}
\end{center}

\begin{abstract}
This note states the weight estimates attached to the six operations.  It
keeps the essential asymmetry visible: pullback and exceptional direct image
preserve upper bounds, while exceptional pullback and ordinary direct image
preserve lower bounds.  The stronger statements for proper or smooth maps
are stated separately.  All categories are ordinary triangulated categories
of strict Weil complexes.
\end{abstract}
\tableofcontents

\section{The categories and their Frobenius structures}

Let \(k=\Fq\), let \(\ell\ne\operatorname{char}k\), and fix
\(\bar k=\Fbar\), \(\Qlb\), and an abstract embedding
\(\iota:\Qlb\hookrightarrow\mathbf C\).  The arithmetic Galois element is
\(a\mapsto a^q\), and the geometric Galois element is its inverse.  If
\(X_0\) is of finite type over \(k\), write \(X=X_0\times_k\bar k\) and
\[
 X^{(q)}=X\times_{\bar k,a\mapsto a^q}\bar k.
\]
The relative \(q\)-Frobenius \(X\to X^{(q)}\), followed by the canonical
descent identification \(X^{(q)}\simeq X\), is the scheme Frobenius
\[
                         \Phi_X:X\longrightarrow X.
\]
It is a universal homeomorphism, not in general an automorphism.  Its
pullback is an exact equivalence on constructible pro-etale sheaves.  Under
the descent comparison, this pullback corresponds to the geometric Galois
action; the fixed geometric points of \(\Phi_X^n\) are
\(X_0(\mathbf F_{q^n})\).

For a finite extension \(E/\Ql\), with valuation ring \(\mathcal O_E\) and
uniformizer \(\varpi_E\), the coefficient sheaf is
\[
 E_X=\left(\varprojlim_n
    (\mathcal O_E/\varpi_E^n)_X\right)[1/\varpi_E],
\]
where the finite-level terms are constant sheaves on the pro-etale site.
Let \(\mathsf{Mod}(X,E)\) be the abelian category of all sheaves of
\(E_X\)-modules and define
\[
 D(X,E)=K(\mathsf{Mod}(X,E))
       [\text{quasi-isomorphisms}^{-1}],
\]
the ordinary derived category obtained from the homotopy category of
unbounded cochain complexes.  An \(E_X\)-sheaf is lisse of rank \(r\) if it is pro-etale locally
\(E_X^{\oplus r}\), and it is constructible if it is lisse of finite rank
on every member of a finite locally closed stratification.  The category
\(\Dbc(X,E)\) consists of complexes with constructible cohomology sheaves,
zero outside a finite interval.

For a finite extension \(E/\Ql\) inside \(\Qlb\), let
\(\mathsf{Cons}_{\mathrm W}(X_0,E)\) be the abelian category whose objects
are constructible \(E\)-sheaves \(\mathcal F\) on \(X_{\mathrm{pro\acute
et}}\) equipped with an isomorphism
\(\Phi_X^*\mathcal F\simeq\mathcal F\), and whose morphisms commute with that
isomorphism.  Define
\(K^b(\mathcal A)\) to be the homotopy category of bounded cochain complexes
in an abelian category \(\mathcal A\), with chain maps modulo chain homotopy.
Then define
\[
 \WeilD(X_0,E)
 =K^b(\mathsf{Cons}_{\mathrm W}(X_0,E))
       [\text{equivariant quasi-isomorphisms}^{-1}].
\]
Thus \(\WeilD(X_0,E)\) is the ordinary Verdier localization of strict
equivariant complexes.  It is not defined as the naive category of pairs in
\(\Dbc(X,E)\).  The exact forgetful functor is
\[
\operatorname{forg}_X:\WeilD(X_0,E)\longrightarrow\Dbc(X,E).
\]

Whenever \(D_X\), \(Rf^!\), or a lower weight predicate is used below, the
schemes involved are assumed separated over \(k\), so their structure
morphisms are compactifiable and Verdier duality is defined on the displayed
bounded constructible categories.

For \(E\subset M\), scalar extension sends \((K,\varphi)\) to
\((K\otimes_EM,\varphi_M)\), with
\[
 \Phi_X^*(K\otimes_EM)\xrightarrow{\sim}
 \Phi_X^*K\otimes_EM\xrightarrow{\varphi\otimes1_M}K\otimes_EM.
\]
This is an exact functor and preserves bounded constructibility.

For \(\Qlb\)-coefficients, take the filtered pair category over finite
extensions \(E/\Ql\), with
\[
 \operatorname{Hom}((E,K),(E',L))
 =\varinjlim_{M\supset E,E'}
 \operatorname{Hom}_{\WeilD(X_0,M)}
       (K\otimes_EM,L\otimes_{E'}M),
\]
form shifts and distinguished triangles at a common finite stage, and then
take its Karoubi envelope.  Denote the resulting idempotent-complete ordinary
triangulated category by
\(\WeilD(X_0,\Qlb)_{\mathrm{2-colim}}\).  In particular,
every object and every operation appearing below has a representative over
one finite coefficient extension of \(\Ql\).
A triangle in the Karoubi envelope is distinguished when it is isomorphic
to a direct summand of a distinguished triangle represented at a common
finite stage.

\section{Equivariant forms of the six operations}

Let \(f:X_0\to Y_0\) be a morphism of finite-type \(k\)-schemes and let
\(f:X\to Y\) also denote its base change to \(\bar k\).  For \(Rf_!\) assume
that \(f\) is separated of finite type; it is the compactification-defined
exceptional direct image.  The usual constructibility and boundedness
theorems give exact functors on \(\Dbc\)
\[
\begin{aligned}
 Lf^*&:\Dbc(Y,E)\longrightarrow\Dbc(X,E),&
 Rf_*&:\Dbc(X,E)\longrightarrow\Dbc(Y,E),\\
 Rf_!&:\Dbc(X,E)\longrightarrow\Dbc(Y,E),&
 Rf^!&:\Dbc(Y,E)\longrightarrow\Dbc(X,E).
\end{aligned}
\]
The constructibility and finite cohomological-dimension theorems assert that
all four displayed functors have exactly these bounded constructible sources
and targets; \(Rf^!\) is the right adjoint of \(Rf_!\).

\begin{definition}[Induced Weil structure]
Form the Cartesian square
\[
\begin{CD}
 X'=X\times_{Y,\Phi_Y}Y @>{g}>> X\\
 @V{f'}VV @VV{f}V\\
 Y @>{\Phi_Y}>> Y .
\end{CD}
\]
Naturality of scheme Frobenius gives
\(h:X\to X'\) with \(g\circ h=\Phi_X\) and \(f'\circ h=f\).
The morphism \(h\) is a universal homeomorphism.  Compact-support base
change for the displayed Cartesian square, followed by the equivalence
\(h^*\), gives the Frobenius-exchange isomorphism
\[
 \Phi_Y^*Rf_!K\xrightarrow{\sim}Rf'_!g^*K
       \xrightarrow{\sim}Rf_!\Phi_X^*K.
\]
For a Weil complex \((K,\varphi_K)\) on \(X_0\), the Weil structure on
\(Rf_!K\) is characterized on the underlying derived category by
\[
 \Phi_Y^*Rf_!K\xrightarrow{\sim}Rf_!\Phi_X^*K
       \xrightarrow{Rf_!(\varphi_K)}Rf_!K.
\]
The same Cartesian construction, using topological invariance, gives the
exchange for \(Rf_*\) and hence its Frobenius structure.  For a Weil complex
\((L,\varphi_L)\) on \(Y_0\), the structures on \(Lf^*L\) and \(Rf^!L\)
are characterized by
\[
 \Phi_X^*Lf^*L\xrightarrow{\sim}Lf^*\Phi_Y^*L
       \xrightarrow{Lf^*(\varphi_L)}Lf^*L,
\]
\[
 \Phi_X^*Rf^!L\xrightarrow{\sim}Rf^!\Phi_Y^*L
       \xrightarrow{Rf^!(\varphi_L)}Rf^!L.
\]
\end{definition}

\begin{theorem}[Strict equivariant realization]
The preceding derived Frobenius structures are realized by exact functors
between the strict derived Weil categories,
\[
 Lf^*_{\mathrm W},\quad Rf_{*,\mathrm W},\quad
 Rf_{!,\mathrm W},\quad Rf^!_{\mathrm W},
\]
whose images under the forgetful functors are naturally isomorphic to the
ordinary six-operation functors.  These lifts commute with finite scalar
extension and hence define stagewise functors on
\(\WeilD(-,\Qlb)_{\mathrm{2-colim}}\).  Their unit, composition, base-change,
projection-formula, and Verdier-duality isomorphisms are Frobenius
equivariant and satisfy the usual coherence identities.
\end{theorem}

For readability the subscript \({\mathrm W}\) is suppressed below; every
functor is understood in this strict equivariant sense.

\begin{definition}[Tensor product, internal Hom, and Verdier dual]
On \(D(X,E)\), derived tensor product is computed from a \(K\)-flat
replacement, and \(R\mathcal Hom_E\) is computed from a \(K\)-injective
replacement in its second variable.  For Weil complexes
\((K,\varphi_K)\) and \((L,\varphi_L)\), the tensor Frobenius is
\(\varphi_K\otimes\varphi_L\), while the internal-Hom Frobenius is
conjugation:
\[
 \Phi_X^*R\mathcal Hom_E(K,L)
 \xrightarrow{\sim}R\mathcal Hom_E(\Phi_X^*K,\Phi_X^*L)
 \xrightarrow{R\mathcal Hom_E(\varphi_K^{-1},\varphi_L)}
 R\mathcal Hom_E(K,L).
\]
These constructions restrict to exact bifunctors
\[
 \begin{aligned}
 -\otimes_E^L- &: \WeilD(X_0,E)\times\WeilD(X_0,E)
                    \longrightarrow\WeilD(X_0,E),\\
 \RHom_E(-,-) &: \WeilD(X_0,E)^{\mathrm{op}}\times\WeilD(X_0,E)
                    \longrightarrow\WeilD(X_0,E).
 \end{aligned}
\]
Define the strict linear dual by
\[
                         K^\vee=\RHom_E(K,E_X).
\]

Put \(S_0=\operatorname{Spec}k\), \(S=\operatorname{Spec}\bar k\), and
let \(a_X:X\to S\) be the structure morphism.  Set
\[
 \omega_X=Ra_X^!E_S,
 \qquad D_XK=R\mathcal Hom_E(K,\omega_X).
\]
Give the constant object \(E_S\) its identity Weil structure.  Exceptional
Frobenius exchange gives
\[
 \varphi_\omega:\Phi_X^*\omega_X
 \xrightarrow{\sim}Ra_X^!\Phi_S^*E_S
 \xrightarrow{\sim}Ra_X^!E_S=\omega_X.
\]
The Weil
structure on the dual is
\[
 \Phi_X^*D_XK\xrightarrow{\sim}D_X(\Phi_X^*K)
       \xrightarrow{(D_X\varphi_K)^{-1}}D_XK.
\]
Thus on a point the dual Frobenius is the inverse transpose.
\end{definition}

\begin{theorem}[Strict monoidal and duality realization]
The preceding formulas are induced by exact functors on the strict derived
Weil category.  Their images under the forgetful functor are the ambient
derived tensor product, internal Hom, and Verdier duality.  In particular,
\[
 D_X:\WeilD(X_0,E)^{\mathrm{op}}\xrightarrow{\sim}\WeilD(X_0,E)
\]
is an exact anti-equivalence with \(K\xrightarrow{\sim}D_XD_XK\).  These
functors preserve bounded constructibility and commute with finite scalar
extension, so they act stagewise on the finite-stage \(\Qlb\)-category.
\end{theorem}

\begin{definition}[Strict Tate objects and twists]
On the pro-etale site put
\[
 \mathbf Z_\ell(1)=\varprojlim_n\mu_{\ell^n},
 \qquad E(1)=\mathbf Z_\ell(1)\otimes_{\mathbf Z_\ell}E,
\]
where \(\mu_{\ell^{n+1}}\to\mu_{\ell^n}\) is the \(\ell\)-th-power map.
Functorial pullback of roots of unity gives a strict arrow
\(\Phi_X^*E(1)\xrightarrow{\sim}E(1)\).  Put \(E(0)=E_X\), define
\(E(m)=E(1)^{\otimes_E^Lm}\) for \(m>0\), and define
\(E(-m)=\RHom_E(E(m),E_X)\) for \(m>0\), with the tensor and
internal-Hom Weil arrows.  For every \(m\in\mathbf Z\), the Tate-twist
functor is the exact functor
\[
 (-)(m)=-\otimes_E^LE(m):\WeilD(X_0,E)\longrightarrow\WeilD(X_0,E).
\]
At a closed point \(x\), geometric Frobenius on \(E(m)_{\bar x}\) is
\(Q_x^{-m}\operatorname{id}\).  The objects and functors commute with
finite scalar extension.
\end{definition}

\section{Finite Frobenius modules and scalar weights}

\begin{definition}[Finite \(Q\)-Frobenius module]
Let \(Q>1\) be an integer.  A finite \(Q\)-Frobenius module over \(E\) is a
triple \((Q,V,F)\), where \(V\) is a finite-dimensional \(E\)-vector space
and \(F\in\operatorname{GL}_E(V)\).  Its spectrum is the multiset of roots
of \(\det(T-F\mid V)\) in \(\Qlb\), counted with multiplicity.
\end{definition}

\begin{definition}[Fixed-embedding \(Q\)-weight predicates]
For a finite \(Q\)-Frobenius module \((Q,V,F)\) and \(w\in\mathbf R\), put
\[
 \Pure_{Q,\iota}(V,F;w)
 \Longleftrightarrow
 |\iota(\alpha)|=Q^{w/2}
 \quad\text{for every eigenvalue }\alpha\text{ of }F.
\]
The predicate \(\MixLe_{Q,\iota}(V,F;w)\) means that \(V\) has a finite
\(F\)-stable filtration whose nonzero graded modules are
\(\Pure_{Q,\iota}\) of weights at most \(w\).  The predicate
\(\MixGe_{Q,\iota}(V,F;w)\) is defined with weights at least \(w\).
\end{definition}

\begin{definition}[All-conjugates \(Q\)-weight predicates]
A nonzero \(\alpha\in\Qlb\) is a \(Q\)-Weil number of weight \(w\) if it is
algebraic over \(\mathbf Q\) and
\[
                         |\sigma(\alpha)|=Q^{w/2}
\]
for every embedding \(\sigma:\mathbf Q(\alpha)\hookrightarrow\mathbf C\).
A finite \(Q\)-Frobenius module is
\(\Pure_Q^{\mathrm{all}}(V,F;w)\) when all its eigenvalues are \(Q\)-Weil
numbers of weight \(w\).  Define \(\UniLe_Q(V,F;w)\) by requiring every
eigenvalue \(\alpha\) to be algebraic over \(\mathbf Q\) and to satisfy
\[
                         |\sigma(\alpha)|\le Q^{w/2}
\]
for every embedding \(\sigma:\mathbf Q(\alpha)\hookrightarrow\mathbf C\).
Define \(\UniGe_Q(V,F;w)\) by replacing \(\leq\) with \(\geq\).
\end{definition}

\begin{definition}[Integral \(Q\)-mixedness]
A finite \(Q\)-Frobenius module is integrally
\(\MixLe_{Q,\iota}(V,F;w)\) if it has a finite \(F\)-stable filtration whose
nonzero graded modules are \(\Pure_{Q,\iota}\) of weights
\(u\in\mathbf Z\) with \(u\le w\).  Integral
\(\MixGe_{Q,\iota}(V,F;w)\) is defined with \(u\ge w\).  It is
all-conjugates integrally mixed of weights at most \(w\) when the same
definition uses \(\Pure_Q^{\mathrm{all}}\) of integral weights \(u\leq w\)
on every nonzero graded module.  The lower predicate uses integral weights
\(u\geq w\).
\end{definition}

\begin{proposition}[Exact sequences]
For every exact \(F\)-equivariant sequence of finite \(Q\)-Frobenius modules
\[
 0\longrightarrow(V',F')\longrightarrow(V,F)
   \longrightarrow(V'',F'')\longrightarrow0,
\]
one has
\[
 \det(T-F\mid V)=\det(T-F'\mid V')\det(T-F''\mid V'').
\]
Consequently the fixed-embedding upper and lower predicates, their integral
versions, and \(\UniLe_Q\) are stable under \(F\)-stable subobjects,
quotients, extensions, finite direct sums, and direct summands.
\end{proposition}

\begin{proposition}[Tensor, dual, and Tate weight algebra]
Let \((Q,V,F_V)\) and \((Q,W,F_W)\) be finite \(Q\)-Frobenius modules.
Give \(V\otimes_EW\) the operator \(F_V\otimes F_W\), give \(V^\vee\)
the operator \((F_V^{-1})^\vee\), and give \(V(r)\) the operator
\(Q^{-r}F_V\).  Then
\begin{enumerate}[label=\textup{\(\roman*\)}]
\item \(\Pure_{Q,\iota}(V,F_V;u)\) and
\(\Pure_{Q,\iota}(W,F_W;v)\) imply
\(\Pure_{Q,\iota}(V\otimes W,F_V\otimes F_W;u+v)\).
\item \(\Pure_{Q,\iota}(V,F_V;u)\) implies
\(\Pure_{Q,\iota}(V^\vee,(F_V^{-1})^\vee;-u)\).
\item \(\Pure_{Q,\iota}(V,F_V;u)\) implies
\(\Pure_{Q,\iota}(V(r),Q^{-r}F_V;u-2r)\).
\item The same statements hold with \(\Pure_{Q,\iota}\) replaced by
\(\Pure_Q^{\mathrm{all}}\), and the corresponding upper bounds hold for
mixed modules and for \(\UniLe_Q\).
\end{enumerate}
\end{proposition}

\begin{definition}[Bounded \(Q\)-Frobenius spectral sequence]
A bounded \(Q\)-Frobenius spectral sequence is a spectral sequence
\((E_r^{a,b},d_r)\) of finite \(Q\)-Frobenius modules whose differentials
commute with Frobenius, which is zero outside a finite rectangle, and which
converges to finite \(Q\)-Frobenius modules \(H^n\) with finite
Frobenius-stable filtrations
\[
                         \gr^aH^n\simeq E_\infty^{a,n-a}.
\]
\end{definition}

\begin{proposition}[Spectral-sequence upper bound]
In a bounded \(Q\)-Frobenius spectral sequence, suppose
\[
             \MixLe_{Q,\iota}(E_1^{a,b},F;w+a+b)
\]
for all \(a,b\).  Then
\[
             \MixLe_{Q,\iota}(H^n,F;w+n)
\]
for every \(n\).  The same assertion holds for integral mixedness and for
the all-conjugates predicate \(\UniLe_Q\).
\end{proposition}

\section{Closed-point spectra and weights of complexes}

\begin{definition}[Closed-point Frobenius module]
Let \(x\in|X_0|\) be closed, put
\(d_x=[\kappa(x):k]\) and \(Q_x=q^{d_x}\), and choose a geometric point
\(\bar x\to X\) above \(x\).  Iterating the strict Weil arrow \(d_x\) times
gives the local geometric Frobenius
\[
 F_x:\mathcal F_{\bar x}\xrightarrow{\sim}\mathcal F_{\bar x}.
\]
Thus \((Q_x,\mathcal F_{\bar x},F_x)\) is a finite
\(Q_x\)-Frobenius module.  Its characteristic polynomial is independent of
the choice of \(\bar x\).
\end{definition}

\begin{definition}[Punctual purity and mixed Weil sheaves]
A constructible Weil sheaf \(\mathcal F\) is punctually
\(\iota\)-pure of weight \(w\) if
\[
 \Pure_{Q_x,\iota}(\mathcal F_{\bar x},F_x;w)
\]
for every closed point \(x\).  It is mixed of weights at most \(w\) if,
after one finite extension of \(E\), it has a finite filtration by Weil
subsheaves whose nonzero graded pieces are punctually \(\iota\)-pure of
weights at most \(w\).  Mixedness of weights at least \(w\) is defined with
the reversed inequality.
\end{definition}

\begin{definition}[Integral and all-conjugates mixed Weil sheaves]
A constructible Weil sheaf is integrally mixed of weights at most \(w\) if
the pure weights of the preceding finite filtration belong to
\(\mathbf Z\) and are at most \(w\).  It is punctually all-conjugates pure
of integral weight \(u\) if
\[
 \Pure_{Q_x}^{\mathrm{all}}(\mathcal F_{\bar x},F_x;u)
\]
for every closed point \(x\).  It is all-conjugates integrally mixed of
weights at most \(w\) if it has a finite Weil filtration with punctually
all-conjugates pure graded pieces of integral weights \(u\leq w\).  The
lower condition uses integral weights \(u\geq w\).
\end{definition}

\begin{definition}[Upper, lower, and pure complexes]
For \(K\in\WeilD(X_0,E)\) and \(w\in\mathbf R\), define
\[
 \WtLe(K,w)\Longleftrightarrow
 \mathcal H^m(K)\text{ is mixed of weights at most }w+m
 \text{ for every }m.
\]
Define \(\IntWtLe(K,w)\) by requiring every cohomology sheaf to be
integrally mixed with the same bound.  Define
\[
 \WtGe(K,w)\Longleftrightarrow\WtLe(D_XK,-w),
 \qquad
 \IntWtGe(K,w)\Longleftrightarrow\IntWtLe(D_XK,-w).
\]
Define \(\AllIntWtLe(K,w)\) by requiring \(\mathcal H^m(K)\) to be
all-conjugates integrally mixed of weights at most \(w+m\), and define
\[
 \AllIntWtGe(K,w)\Longleftrightarrow\AllIntWtLe(D_XK,-w).
\]
The complex is pure of weight \(w\) if both \(\WtLe(K,w)\) and
\(\WtGe(K,w)\) hold; it is integrally pure if the two integral predicates
hold.
\end{definition}

\begin{proposition}[Point criterion]
For \(X_0=\Spec k\), an object \(K\in\WeilD(X_0,E)\) satisfies
\(\WtLe(K,w)\) if and only if
\[
 \MixLe_{q,\iota}(\mathcal H^m(K),F;w+m)
\]
for every \(m\).  It satisfies \(\WtGe(K,w)\) if and only if
\[
 \MixGe_{q,\iota}(\mathcal H^m(K),F;w+m)
\]
for every \(m\).  Hence purity of \(K\) of weight \(w\) is equivalent to
\(\Pure_{q,\iota}(\mathcal H^m(K),F;w+m)\) for every \(m\).
\end{proposition}

\section{Exact weight calculus}

\begin{proposition}[Abelian stability]
For every real \(w\), constructible Weil sheaves mixed of weights at most
\(w\) are closed under Weil subobjects, quotients, extensions, finite direct
sums, and direct summands.  The same holds for integral mixedness and for
all-conjugates integral mixedness.  All three conditions are invariant under
finite extension of the coefficient field.
\end{proposition}

\begin{proposition}[Shifts, twists, and Verdier duality]
For \(K\in\WeilD(X_0,E)\) and integers \(r,s\),
\[
 \WtLe(K[r],w+r)\Longleftrightarrow\WtLe(K,w),
 \qquad
 \WtLe(K(s),w-2s)\Longleftrightarrow\WtLe(K,w).
\]
The same equivalences hold for \(\WtGe\), and
\[
 \WtLe(K,w)\Longleftrightarrow\WtGe(D_XK,-w),
 \qquad
 \WtGe(K,w)\Longleftrightarrow\WtLe(D_XK,-w).
\]
All assertions remain valid for the integral predicates.
\end{proposition}

\begin{proposition}[Correct triangle bounds]
Let \(A\to B\to C\to A[1]\) be a distinguished triangle.  Then
\begin{align*}
 \WtLe(A,a),\ \WtLe(C,c)&\Longrightarrow
     \WtLe(B,\max\{a,c\}),\\
 \WtLe(A,a),\ \WtLe(B,b)&\Longrightarrow
     \WtLe(C,\max\{a+1,b\}),\\
 \WtLe(B,b),\ \WtLe(C,c)&\Longrightarrow
     \WtLe(A,\max\{b,c-1\}).
\end{align*}
For the lower predicates one has
\begin{align*}
 \WtGe(A,a),\ \WtGe(C,c)&\Longrightarrow
     \WtGe(B,\min\{a,c\}),\\
 \WtGe(A,a),\ \WtGe(B,b)&\Longrightarrow
     \WtGe(C,\min\{a+1,b\}),\\
 \WtGe(B,b),\ \WtGe(C,c)&\Longrightarrow
     \WtGe(A,\min\{b,c-1\}).
\end{align*}
The six implications remain valid for integral weights.
\end{proposition}

\begin{theorem}[Tensor weight theorem]
For \(K,L\in\WeilD(X_0,E)\),
\[
 \WtLe(K,u),\ \WtLe(L,v)
 \quad\Longrightarrow\quad
 \WtLe(K\otimes_E^LL,u+v).
\]
The same statement holds for integral upper weights.  If the inputs are
pure of weights \(u\) and \(v\), the tensor product is pure of weight
\(u+v\).
\end{theorem}

\begin{theorem}[Internal-Hom weight theorem]
Let \(K,L\in\WeilD(X_0,E)\).  Assume that \(K\) is dualizable, meaning
that the canonical evaluation map
\[
             K^\vee\otimes_E^LL\xrightarrow{\sim}\RHom_E(K,L)
\]
is an isomorphism, and assume \(\WtLe(K^\vee,-u)\).  Then
\[
 \WtLe(L,v)
 \quad\Longrightarrow\quad
 \WtLe(\RHom_E(K,L),v-u).
\]
The same assertion holds for integral weights.
\end{theorem}

\begin{corollary}[Smooth Verdier form of the internal-Hom estimate]
Assume that \(X_0\) is smooth of pure dimension \(d\), so that
\(\omega_X\simeq E(d)[2d]\).  For a dualizable \(K\), the comparison
\[
 K^\vee\simeq D_XK(-d)[-2d]
\]
shows that \(\WtGe(K,u)\) implies \(\WtLe(K^\vee,-u)\).  Hence
\(\WtGe(K,u)\) and \(\WtLe(L,v)\) imply
\(\WtLe(\RHom_E(K,L),v-u)\), with the identical integral statement.
\end{corollary}

\begin{proposition}[Local detection]
Let \(X_0=\bigsqcup_{a=1}^rS_a\) be a finite stratification by locally
closed subschemes, with inclusions \(s_a\).  A constructible Weil sheaf
already equipped with a mixed filtration has weights at most \(w\) if and
only if every \(s_a^*\mathcal F\) has weights at most \(w\).  A complex
\(K\) satisfies \(\WtLe(K,w)\) if and only if the same test holds for every
\(s_a^*\mathcal H^m(K)\) with bound \(w+m\).  The assertions also hold for
integral and all-conjugates integral mixedness.
\end{proposition}

\begin{proposition}[Finite extension of the ground field]
Let \(k_r=\mathbf F_{q^r}\), let \(x\) have degree \(d\) over \(k\), and
put \(g=\gcd(d,r)\).  Over \(k_r\), the point \(x\) splits into \(g\)
closed points.  Each has residue-field cardinality
\[
 Q'=q^{rd/g}=Q_x^{r/g}
\]
and local geometric Frobenius \(F_x^{r/g}\).  Punctual purity, mixedness, integral mixedness, all-conjugates
mixedness, and the upper and lower complex predicates are invariant under
finite extension of the ground field.
\end{proposition}

\section{Elementary functorial weight statements}

\begin{theorem}[Pullback preserves upper weights]
For every \(f_0:X_0\to Y_0\), the exact strict functor
\(Lf^*:\WeilD(Y_0,E)\to\WeilD(X_0,E)\) satisfies
\[
             \WtLe(K,w)\Longrightarrow\WtLe(Lf^*K,w).
\]
It also preserves \(\IntWtLe\) and all-conjugates integral upper weights.
At a closed \(x\in X_0\) over \(y\in Y_0\), the local finite Frobenius
modules use \(Q_x=q^{\deg(x)}\) and \(Q_y=q^{\deg(y)}\), and the strict
pullback comparison identifies the appropriate positive iterates.
\end{theorem}

\begin{theorem}[Exceptional pullback preserves lower weights]
Let \(f_0:X_0\to Y_0\) be separated and of finite type.  The strict
exceptional pullback satisfies
\[
             \WtGe(K,w)\Longrightarrow\WtGe(Rf^!K,w).
\]
The same statement holds for integral lower weights.
\end{theorem}

\begin{theorem}[Finite direct image]
If \(u_0:T_0\to S_0\) is finite, then
\[
 Ru_!=Ru_*=u_*:\WeilD(T_0,E)\longrightarrow\WeilD(S_0,E),
\]
and \(u_*\) is exact on constructible Weil sheaves.  It preserves punctual
purity, upper and lower weights, integral upper and lower weights, and
all-conjugates integral weights.  For a geometric point \(\bar s\to S_0\),
\[
 (u_*\mathcal F)_{\bar s}
 \simeq\bigoplus_{\bar t\in T_{\bar s}}\mathcal F_{\bar t}
\]
with the Frobenius permutation and the local iterate actions.
\end{theorem}

\begin{corollary}[Closed immersions]
For a closed immersion \(i_0:Z_0\hookrightarrow X_0\),
\(i_!=i_*\) is exact, fully faithful, and preserves upper and lower weights
without changing the bound.  The pullback \(Li^*\) preserves upper weights,
 and the exceptional pullback \(Ri^!\) preserves lower weights.
All statements retain their integral and all-conjugates variants.
\end{corollary}

\begin{theorem}[Open immersions and localization]
Let \(j_0:U_0\hookrightarrow X_0\) be open with closed complement
\(i_0:Z_0\hookrightarrow X_0\).  There are Frobenius-equivariant
distinguished triangles
\[
 j_!j^*K\longrightarrow K\longrightarrow i_*i^*K
                   \longrightarrow j_!j^*K[1]
\]
and
\[
 i_*Ri^!K\longrightarrow K\longrightarrow Rj_*j^*K
                   \longrightarrow i_*Ri^!K[1].
\]
The functors \(j^*=j^!\) preserve upper and lower weights, \(j_!\)
preserves upper weights, and \(Rj_*\) preserves lower weights.  The same
assertions hold for integral weights.  Moreover,
\[
 \WtLe(K,w)\Longleftrightarrow
 \WtLe(j^*K,w)\text{ and }\WtLe(i^*K,w),
\]
with an identical equivalence for \(\IntWtLe\).
\end{theorem}

\begin{theorem}[Separated quasi-finite direct image]
Let \(u_0:T_0\to S_0\) be separated and quasi-finite.  Then \(Ru_!\)
preserves upper weights and is concentrated in degree zero on sheaves.  It
preserves integral and all-conjugates integral upper weights.  A finite
ordered stratification of \(T_0\) may be chosen so that each restriction
factors as an open immersion followed by a finite morphism.
\end{theorem}

\begin{definition}[Smooth-purity hypotheses]
For a smooth morphism \(f_0:X_0\to Y_0\) of pure relative dimension \(d\),
the smooth-purity datum is the Frobenius-equivariant natural isomorphism
\[
                    Rf^!(-)\xrightarrow{\sim}Lf^*(-)(d)[2d]
\]
compatible with finite scalar extension, composition, and base change.
\end{definition}

\begin{theorem}[Smooth pullback]
If \(f_0\) is smooth of pure relative dimension \(d\) and the smooth-purity
hypotheses hold, then \(Lf^*\) preserves both upper and lower weights
without changing the bound.  It preserves purity, integral purity, and
all-conjugates integral purity.  The shift by \(2d\) and twist by \(d\) in
\(Rf^!\simeq Lf^*(d)[2d]\) have net weight zero.
\end{theorem}

\begin{corollary}[Regular closed immersion]
Let \(i_0:Z_0\hookrightarrow X_0\) be a regular closed immersion of pure
codimension \(c\), and assume the Frobenius-equivariant absolute-purity
isomorphism
\[
                    Ri^!(-)\xrightarrow{\sim}Li^*(-)(-c)[-2c].
\]
Then \(Li^*\) and \(Ri^!\) preserve both weight inequalities and purity,
including their integral and all-conjugates forms.
\end{corollary}

\section{Integral stratification and bounded devissage}

\begin{definition}[Ordered support stratification]
An ordered support stratification of a noetherian \(X_0\) is a chain
\[
 \varnothing=X_{-1}\subset X'_0\subset\cdots\subset X'_r=X_0
\]
of closed reduced subschemes such that
\(S_a=X'_a\setminus X'_{a-1}\) is smooth and locally closed.  Write
\(s_a:S_a\hookrightarrow X_0\) for the inclusion.
\end{definition}

\begin{theorem}[Pure-lisse stratification]
Let \(\mathcal F\) be an integrally mixed constructible Weil sheaf of
weights at most \(w\).  After a finite extension of \(E\), there are an
ordered support stratification and a finite Frobenius-stable filtration of
\(\mathcal F\) whose successive quotients have the form
\[
                         s_{a,!}\mathcal L_a,
\]
where \(\mathcal L_a\) is lisse and punctually pure of an integral weight
\(u_a\le w\) on \(S_a\).  A common refinement exists for any finite family
of sheaves, morphisms, kernels, cokernels, and filtrations.
\end{theorem}

\begin{lemma}[One-step source localization]
Let \(j:U_0\hookrightarrow X_0\) be open with closed complement \(i\), and
let \(\mathcal F\) be a constructible Weil sheaf.  If
\[
 R^n(fj)_!j^*\mathcal F
 \quad\text{and}\quad
 R^n(fi)_!i^*\mathcal F
\]
are integrally mixed of weights at most \(w+n\) for every \(n\), then so is
\(R^nf_!\mathcal F\).  The assertion includes the Frobenius-equivariant
long exact sequence obtained from the upper localization triangle.
\end{lemma}

\begin{theorem}[Sheaf devissage]
Fix \(f_0:X_0\to Y_0\).  Suppose that for every smooth locally closed
\(s:S_0\hookrightarrow X_0\) and every lisse punctually pure Weil sheaf
\(\mathcal L\) of integral weight \(u\),
\[
              R^m(fs)_!\mathcal L
 \text{ is integrally mixed of weights at most }u+m.
\]
Then, for every integrally mixed constructible Weil sheaf \(\mathcal F\)
of weights at most \(w\),
\[
              R^mf_!\mathcal F
 \text{ is integrally mixed of weights at most }w+m.
\]
\end{theorem}

\begin{definition}[Postnikov data]
Let \(K\in\WeilD(X_0,E)\), and choose \(a_-\le a_+\) with
\(\mathcal H^a(K)=0\) outside \([a_-,a_+]\).  Put
\(K_a=\tau_{\le a}K\), \(K_{a_--1}=0\), and \(K_{a_+}=K\).  The standard
Frobenius-equivariant Postnikov triangles are
\[
 K_{a-1}\longrightarrow K_a\longrightarrow
 \mathcal H^a(K)[-a]\longrightarrow K_{a-1}[1].
\]
Their length is \(a_+-a_-+1\).
\end{definition}

\begin{proposition}[Postnikov output filtration]
Assume \(Rf_!\) has cohomological amplitude \([b_-,b_+]\) on the finitely
many cohomology sheaves of \(K\).  For every \(n\), the images of the
Postnikov triangles give a finite Frobenius-stable filtration of
\(\mathcal H^n(Rf_!K)\) whose \(a\)-th graded object is a subquotient of
\[
                       R^{n-a}f_!\mathcal H^a(K).
\]
The graded object is zero unless \(a\in[a_-,a_+]\) and
\(n-a\in[b_-,b_+]\).
\end{proposition}

\begin{theorem}[Bounded-complex devissage]
Suppose the sheaf form of the integral upper-weight theorem holds for
\(f_0\).  Then, for every \(K\in\WeilD(X_0,E)\) and \(w\in\mathbf R\),
\[
               \IntWtLe(K,w)\Longrightarrow
               \IntWtLe(Rf_!K,w).
\]
Equivalently, for every \(n\), the finite filtration in the preceding
proposition has graded subquotients integrally mixed of weights at most
\(w+n\).
\end{theorem}

\section{The integral affine-line theorem}

Let \(A_0=\mathbf A^1_k\), let \(j:U_0\hookrightarrow A_0\) be dense, and
let \(A'_0\) be the dual affine line.  Write \(p,p'\) for the projections
from \(A\times A'\), and let
\(\mu:A\times A'\to\mathbf A^1_{\bar k}\) be multiplication.

\begin{definition}[Artin--Schreier sheaf]
After replacing \(E\) by a finite extension containing \(\mu_p\), fix a
nontrivial additive character
\(\psi:\mathbf F_p\to E^\times\), and put
\(\psi_k=\psi\circ\operatorname{Tr}_{k/\mathbf F_p}\).  Let
\(\pi:\mathbf A^1\to\mathbf A^1\) be \(z\mapsto z^p-z\).  Translation by
\(a\in\mathbf F_p\) acts on \(\pi_*E\), and the idempotent
\[
 e_\psi=p^{-1}\sum_{a\in\mathbf F_p}\psi(a)^{-1}[a]
\]
defines the rank-one Artin--Schreier summand \(\mathcal L_\psi\).  Its
strict Weil structure is the one induced by the geometric Frobenius action
 on the cover and the character
\(\psi_k\circ\operatorname{Tr}_{\mathbf F_{q^n}/k}\) over every finite
extension \(\mathbf F_{q^n}/k\).  Further finite scalar extension does not
change this strict Weil sheaf or any of its weight predicates.
\end{definition}

\begin{theorem}[Artin--Schreier package]
The morphism \(\pi\) is a finite etale
\((\mathbf F_p,+)\)-torsor.  The projectors \(e_\psi\) are pairwise
orthogonal, their sum is the identity, and
\[
                       \pi_*E\simeq\bigoplus_\psi\mathcal L_\psi
\]
as constructible Weil sheaves.  Every \(\mathcal L_\psi\) is lisse of rank
one, its formation commutes with finite extension of coefficients and of
the ground field, and
\(\mathcal L_\psi^\vee\simeq\mathcal L_{\psi^{-1}}\).
\end{theorem}

\begin{definition}[Fourier--Deligne transform]
For \(K\in\WeilD(A_0,E)\), define
\[
 \operatorname{FT}_\psi(K)
   =Rp'_!\bigl(p^*K\otimes_E^L\mu^*\mathcal L_\psi\bigr)[1]
   \in\WeilD(A'_0,E).
\]
The strict Weil arrow is the composite of the pullback, tensor, kernel, and
compact-support Frobenius exchanges.  Thus
\[
 \operatorname{FT}_\psi:\WeilD(A_0,E)\longrightarrow
                              \WeilD(A'_0,E)
\]
is an exact functor of ordinary triangulated categories.
\end{definition}

\begin{theorem}[Fourier constructibility and inversion]
The Fourier--Deligne transform preserves bounded constructibility and
commutes with finite extension of coefficients and of the ground field.  If
\(a:A\to A\) is \(x\mapsto-x\), there is a Frobenius-equivariant natural
isomorphism
\[
 \operatorname{FT}_{\psi^{-1}}\operatorname{FT}_\psi(K)
                    \xrightarrow{\sim}a^*K(-1).
\]
The two inversion isomorphisms satisfy the triangle identities.
\end{theorem}

\begin{definition}[Type-3 affine-line sheaf]
A type-3 sheaf of fixed-embedding weight \(w\) is a lisse Weil sheaf
\(\mathcal L\) on \(U_0\) satisfying the following conditions.
\begin{enumerate}[label=\textup{\(\roman*\)}]
\item \(\mathcal L\) is irreducible in
\(\mathsf{Cons}_{\mathrm W}(U_0,E)\) and punctually
\(\iota\)-pure of weight \(w\).
\item The geometric inertia group at
\(\infty\in\mathbf P^1_k\) acts trivially.
\item After scalar extension to \(\Qlb\), the geometric local system has no
simple subquotient which, for some finite extension \(k'/k\) and some
\(a\in k'\), becomes the Artin--Schreier sheaf
\(\mathcal L_{\psi_k\circ\operatorname{Tr}_{k'/k}}(ax)\).  The value
\(a=0\) is included, so geometrically constant constituents are excluded.
\end{enumerate}
It is an all-conjugates type-3 sheaf of integral weight
\(w\in\mathbf Z\) if the punctual condition in \(i\) is
all-conjugates purity relative to every \(Q_x=q^{\deg(x)}\).
\end{definition}

\begin{definition}[Local Fourier module]
Let \(\mathcal G=\operatorname{FT}_\psi(j_*\mathcal L[1])\), and suppose
its restriction to a punctured strict henselian neighborhood of
\(0'\in A'_0\) is lisse and tamely ramified.  Its local Fourier module is
the finite \(q\)-Frobenius module \((q,V',F)\) given by the geometric
generic stalk, together with the nilpotent tame logarithm
\[
               N:V'\longrightarrow V'(-1),
 \qquad FNF^{-1}=q^{-1}N.
\]
The monodromy filtration centered at zero is the unique finite increasing
filtration \(M_\bullet V'\) for which
\(N(M_a)\subset M_{a-2}(-1)\) and
\[
 N^r:\gr_r^MV'\xrightarrow{\sim}\gr_{-r}^MV'(-r)
 \qquad(r\ge0).
\]
\end{definition}

\begin{theorem}[Type-3 Fourier shape and stationary phase]
For a type-3 sheaf \(\mathcal L\), the normalized transform of
\(j_*\mathcal L[1]\) is a middle extension of a lisse sheaf on the
punctured dual line.  Its local Fourier module at \(0'\) is tame and there
is a Frobenius-equivariant isomorphism
\[
              H_c^1(A,j_*\mathcal L)\xrightarrow{\sim}\ker N.
\]
Moreover \(H_c^r(A,j_*\mathcal L)=0\) for \(r\ne1\).
\end{theorem}

\begin{theorem}[Integral local-Fourier filtration]
Let \(\mathcal L\) be type-3 of integral fixed-embedding weight
\(w\in\mathbf Z\), and let \((q,V',F,N)\) be its local Fourier module.
Then \(N^2=0\), and
\[
 0\subset M_{-1}=\operatorname{im}N
 \subset M_0=\ker N\subset M_1=V'.
\]
After one finite extension of \(E\),
\begin{align*}
 \Pure_{q,\iota}(\gr_{-1}^MV',F;w),\qquad
 \Pure_{q,\iota}(\gr_0^MV',F;w+1),\qquad
 \Pure_{q,\iota}(\gr_1^MV',F;w+2),
\end{align*}
and
\(N:\gr_1^MV'\xrightarrow{\sim}\gr_{-1}^MV'(-1)\).
The stationary-phase isomorphism identifies the two-step filtration
\[
 0\subset\operatorname{im}N\subset\ker N
\]
with a Frobenius-stable filtration of
\(H_c^1(A,j_*\mathcal L)\).
\end{theorem}

\begin{theorem}[Type-3 fixed-embedding endpoint]
If \(\mathcal L\) is type-3 of integral fixed-embedding weight
\(w\in\mathbf Z\), then
\[
 \MixLe_{q,\iota}
 \bigl(H_c^1(A,j_*\mathcal L),F;w+1\bigr),
\]
and the other compact-support groups vanish.  The assertion is unchanged by
finite extension of \(E\) or \(k\).
\end{theorem}

\begin{theorem}[Type-3 all-conjugates endpoint]
If \(\mathcal L\) is all-conjugates type-3 of integral weight
\(w\in\mathbf Z\), then
\[
 \UniLe_q\bigl(H_c^1(A,j_*\mathcal L),F;w+1\bigr).
\]
If the integral local-Fourier filtration is all-conjugates pure on its
three graded modules, then
\(H_c^1(A,j_*\mathcal L)\) is all-conjugates integrally mixed of weights at
most \(w+1\).
\end{theorem}

\begin{definition}[Additional real-weight hypotheses]
For a type-3 sheaf of real fixed-embedding weight \(w\), a
boundary-conjugate companion is a type-3 sheaf whose closed-point and
boundary-invariant characteristic polynomials are carried by \(\iota\) to
the complex conjugates of those of \(\mathcal L\).  The Fourier-output
purity hypothesis is the assertion that the lisse restriction of
\(\operatorname{FT}_\psi(j_*\mathcal L[1])\) is punctually
\(\iota\)-pure of the predicted weight \(w+1\).
\end{definition}

\begin{theorem}[Conditional real-weight affine-line endpoint]
Let \(\mathcal L\) be type-3 of real fixed-embedding weight \(w\).  If a
boundary-conjugate companion exists or the Fourier-output purity hypothesis
holds, then
\[
 \MixLe_{q,\iota}
 \bigl(H_c^1(A,j_*\mathcal L),F;w+1\bigr),
\]
and \(H_c^r(A,j_*\mathcal L)=0\) for \(r\ne1\).
\end{theorem}

\section{From the affine line to projective curves}

Let \(C_0\) be a smooth projective geometrically connected curve,
let \(j:U_0\hookrightarrow C_0\) be dense, and let \(\mathcal L\) be an
irreducible lisse Weil sheaf on \(U_0\).

\begin{definition}[Good-infinity finite projection]
A good-infinity projection for \((C_0,U_0,\mathcal L)\) consists of a finite
morphism
\[
                    h:C_0\longrightarrow\mathbf P^1_k
\]
and a dense open \(W_0\subset U_0\) such that \(h\) is etale on \(W_0\),
the complement of \(W_0\) maps into \(\mathbf A^1_k\), and every
irreducible nonconstant constituent of
\(h_*(\mathcal L|_{W_0})\) is unramified at infinity.  A finite extension of
the ground field is allowed in the datum.
\end{definition}

\begin{theorem}[Existence of a good-infinity projection]
After a finite extension of \(k\), every triple
\((C_0,U_0,\mathcal L)\) has a good-infinity finite projection.  The finite
pushforward is exact, commutes with the relevant local Frobenius iterates,
and preserves fixed-embedding and all-conjugates integral purity.
\end{theorem}

\begin{proposition}[Finite-pushforward constituent decomposition]
For a good-infinity projection, after a finite extension of \(E\), the
semisimplification of \(h_*(\mathcal L|_{W_0})\) is a finite direct sum of
geometrically constant constituents and middle extensions of irreducible
lisse sheaves on dense opens of \(\mathbf A^1_k\).  Every nonconstant
constituent is type-3 of the same integral weight as \(\mathcal L\).  If
\(\mathcal L\) is all-conjugates pure, the constituents are all-conjugates
pure.
\end{proposition}

\begin{proposition}[Affine--projective localization]
Let \(a:\mathbf A^1\hookrightarrow\mathbf P^1\) have complement
\(i_\infty\).  For every constructible Weil sheaf \(\mathcal G\) on
\(\mathbf P^1\), the localization triangle gives a
Frobenius-equivariant exact sequence
\[
 \cdots\to H_c^r(\mathbf A^1,a^*\mathcal G)
 \to H^r(\mathbf P^1,\mathcal G)
 \to H^r(\infty,i_\infty^*\mathcal G)
 \to\cdots.
\]
For a geometrically constant finite \(q\)-Frobenius module \((q,M,F)\),
\[
 H_c^r(\mathbf A^1,M)=0\ (r\ne2),
 \qquad H_c^2(\mathbf A^1,M)\simeq M(-1).
\]
\end{proposition}

\begin{proposition}[Punctual and constant constituents]
Let \((q,M,F)\) be a finite \(q\)-Frobenius module satisfying
\(\Pure_{q,\iota}(M,F;w)\).  The constant Weil sheaf attached to \(M\)
satisfies
\[
 \Pure_{q,\iota}(H^0(\mathbf P^1,M),F;w),
 \qquad
 \Pure_{q,\iota}(H^2(\mathbf P^1,M),F;w+2),
\]
and has zero \(H^1\).  A punctual sheaf pure of weight \(w\) has cohomology
only in degree zero, with the same weight.  The statements retain the
all-conjugates predicates \(\Pure_q^{\mathrm{all}}\).
\end{proposition}

\begin{theorem}[Projective-curve upper bound]
Let \(\mathcal L\) be punctually \(\iota\)-pure of integral weight
\(w\in\mathbf Z\).  Then, for every \(r\),
\[
 \MixLe_{q,\iota}
 \bigl(H^r(C,j_*\mathcal L),F;w+r\bigr).
\]
The same upper bound holds after every finite extension of \(k\) and \(E\).
If \(\mathcal L\) is all-conjugates pure, then
\[
 \UniLe_q\bigl(H^r(C,j_*\mathcal L),F;w+r\bigr).
\]
\end{theorem}

\begin{theorem}[Dual completion on a curve]
Assume Poincare--Verdier duality for \(C_0\).  For a lisse sheaf
\(\mathcal L\) punctually pure of integral weight \(w\), the dual sheaf
\(\mathcal L^\vee\) is punctually pure of weight \(-w\), and there is a
perfect Frobenius-equivariant pairing
\[
 H^r(C,j_*\mathcal L)\otimes_E
 H^{2-r}(C,j_*\mathcal L^\vee)
 \longrightarrow E(-1).
\]
The projective-curve upper bounds for \(\mathcal L\) and
\(\mathcal L^\vee\) imply
\[
 \Pure_{q,\iota}
 \bigl(H^r(C,j_*\mathcal L),F;w+r\bigr).
\]
The all-conjugates version holds when the two input sheaves are
all-conjugates pure.
\end{theorem}

\begin{corollary}[Constructible curve estimate]
Let \(a:C_0\to\Spec k\) be a separated curve and let
\(K\in\WeilD(C_0,E)\).  Then
\[
 \IntWtLe(K,w)\Longrightarrow\IntWtLe(Ra_!K,w).
\]
For a sheaf \(\mathcal F\) integrally mixed of weights at most \(w\), this
means that the finite \(q\)-Frobenius module
\((q,H_c^r(C,\mathcal F),F)\) is integrally
\(\MixLe_{q,\iota}\) of weights at most \(w+r\).
\end{corollary}

\section{Controlled relative curves}

\begin{definition}[Relative-curve compactification datum]
Let \(S_0\) be connected and normal.  A relative-curve compactification
datum consists of
\[
 X_0\xrightarrow{j}\overline X_0\xrightarrow{\bar f}S_0,
 \qquad D_0\xrightarrow{i}\overline X_0,
\]
where \(j\) is open, \(D_0\) is the reduced closed complement of \(X_0\),
\(i\) is its closed immersion, and \(\bar f\) is smooth and proper of pure relative
dimension one with geometrically connected fibers, and
\(D_0\to S_0\) is finite etale.  The original morphism is
\(f=\bar f j\).
\end{definition}

\begin{definition}[Tame coefficient along the boundary]
For a relative-curve compactification datum, a lisse Weil sheaf
\(\mathcal L\) on \(X_0\) is tame along \(D_0\) if wild inertia acts
trivially on every geometric generic stalk at the boundary.  The residual
tame action is a finite semisimple action followed by a unipotent action
whose logarithm is a nilpotent arrow
\[
                         N:V\longrightarrow V(-1).
\]
\end{definition}

\begin{theorem}[Generic controlled compactification]
Let \(h_0:Z_0\to S_0\) be dominant and separated of finite type between
integral schemes, with generic fiber of dimension one, and let
\(\mathcal L\) be lisse and punctually pure of integral weight on a dense
open of \(Z_0\).  After shrinking \(S_0\), a finite radicial base change,
and a finite etale cover, there are a dense open of \(Z_0\), a
relative-curve compactification datum, and a pullback of \(\mathcal L\)
which is tame along the boundary.
\end{theorem}

\begin{theorem}[Trace retraction for the controlling cover]
Let \(u:T'_0\to T_0\) be the finite etale cover in a generic controlled
compactification, of degree \(N\).  The unit and trace are morphisms of Weil
sheaves and satisfy
\[
 E_{T_0}\xrightarrow{\mathrm{unit}}u_*E_{T'_0}
 \xrightarrow{N^{-1}\mathrm{Tr}}E_{T_0}
\]
has composite \(\operatorname{id}_{E_{T_0}}\).
After tensoring with a lisse coefficient and applying compactly supported
direct image, the original output is a Frobenius-stable direct summand of
the output on the cover.
\end{theorem}

\begin{definition}[Monodromy-pure tame local datum]
Let \(s\) be a boundary point with residue-field cardinality \(Q_s\).  A
monodromy-pure tame local datum of integral weight \(w\) is a finite
\(Q_s\)-Frobenius module \((Q_s,V,F_s)\), a nilpotent
\(N:V\to V(-1)\), and its monodromy filtration \(M_\bullet\), such that
\[
 \Pure_{Q_s,\iota}^{\vphantom{A}}
       (\gr_a^MV,F_s;w+a)
\]
for every nonzero graded module.  In the all-conjugates version, this
predicate is replaced by \(\Pure_{Q_s}^{\mathrm{all}}\).
\end{definition}

\begin{theorem}[Local tame monodromy theorem]
Let \(\mathcal L\) be lisse, punctually pure of integral weight \(w\), and
tame on a punctured smooth curve.  At every missing point, its finite
semisimple inertia eigenspaces, the nilpotent logarithm, and the monodromy
filtration form monodromy-pure tame local data of weight \(w\).  Formation
of these data commutes with finite extension of coefficients and residue
fields.
\end{theorem}

\begin{proposition}[Invariant upper bound]
For a monodromy-pure tame local datum of weight \(w\),
\(\ker N\) and the full inertia-invariant subspace
\[
                    V^I=V^{I_{\mathrm{fin}}}\cap\ker N
\]
are integrally
\(\MixLe_{Q_s,\iota}\) of weights at most \(w\).  In the all-conjugates
case, \(V^I\) is all-conjugates integrally mixed of weights at most \(w\).
\end{proposition}

\begin{theorem}[Spreading of monodromy data]
After shrinking \(S_0\), the finite inertia action, \(N\), the monodromy
filtration, its graded modules, \(\ker N\), and \(V^I\) spread to
constructible Weil sheaves along \(D_0\).  Their formation commutes with
geometric fibers.  Over a closed point \(s\), the comparison intertwines
the local geometric Frobenius after the
\([\kappa(s):k]\)-fold descent iterate and uses
\(Q_s=q^{[\kappa(s):k]}\).
\end{theorem}

\begin{theorem}[Relative tame boundary]
Let \((X_0,\overline X_0,S_0,D_0)\) be a relative-curve compactification
datum and let \(\mathcal L\) be lisse, punctually pure of integral weight
\(w\), and tame along \(D_0\).  Then
\[
                         B=i^*j_*\mathcal L
\]
is constructible and integrally mixed of weights at most \(w\).  Its stalk
at a geometric boundary point is the inertia-invariant module \(V^I\).
The boundary filtration and every graded quotient commute with geometric
fibers and with the stated closed-point Frobenius iterates.
\end{theorem}

\begin{theorem}[Relative boundary localization]
There is a short exact sequence of constructible Weil sheaves
\[
 0\longrightarrow j_!\mathcal L\longrightarrow j_*\mathcal L
   \longrightarrow i_*B\longrightarrow0.
\]
Put
\[
 A^r=R^rf_!\mathcal L,
 \quad P^r=R^r\bar f_*j_*\mathcal L,
 \quad E_B=(\bar f i)_*B.
\]
Then \(A^r=P^r=0\) outside \(0\le r\le2\), and the induced
Frobenius-equivariant long exact sequence contains
\[
 0\to A^0\to P^0\to E_B\to A^1\to P^1\to0,
 \qquad A^2\xrightarrow{\sim}P^2.
\]
All displayed terms and arrows commute with geometric-fiber base change.
\end{theorem}

\begin{theorem}[Projective-fiber comparison]
After shrinking \(S_0\), every \(P^r\) is lisse, and proper base change
gives
\[
 (P^r)_{\bar s}\xrightarrow{\sim}
 H^r(\overline X_{\bar s},j_{\bar s,*}\mathcal L_{\bar s}).
\]
At a closed \(s\), this is an isomorphism of finite
\(Q_s\)-Frobenius modules, where \(Q_s=q^{[\kappa(s):k]}\).  The right-hand
side is \(\Pure_{Q_s,\iota}\) of weight \(w+r\).
\end{theorem}

\begin{theorem}[Controlled relative-curve estimate]
Under the relative tame-boundary hypotheses,
\[
             R^rf_!\mathcal L
 \text{ is integrally mixed of weights at most }w+r
 \qquad(0\le r\le2).
\]
The all-conjugates conclusion holds when the coefficient, the local
monodromy data, and the projective-fiber comparisons have their
all-conjugates forms.
\end{theorem}

\begin{corollary}[Descent from the controlling cover]
In a generic controlled compactification, the trace retraction, finite
base-change comparison, and radicial invariance identify every
\(R^rh_!\mathcal L\) as a Frobenius-stable direct summand of an output
satisfying the controlled relative-curve estimate.  Therefore
\(R^rh_!\mathcal L\) is integrally mixed of weights at most \(w+r\) on the
chosen dense open of the target.
\end{corollary}

\section{Finite reductions and relative affine lines}

\begin{definition}[Support complexity]
For a nonempty constructible support \(Z\) in a finite-type scheme, define
\[
              \mu(Z)=(\dim Z,c(Z))\in\mathbf N\times\mathbf N,
\]
where \(c(Z)\) is the number of irreducible components of maximal dimension.
Use lexicographic order.  The empty support is assigned the least value.
\end{definition}

\begin{theorem}[Finite support devissage]
Let \(\mathcal P(Z,\mathcal F)\) be a property of constructible Weil
sheaves supported on \(Z\).  Suppose it is stable under finite scalar
extension, subquotients, extensions, and direct summands; suppose it glues
across every open--closed decomposition; and suppose it holds for every
pure lisse constituent of integral weight on a dense smooth open of every
irreducible support.  Then it holds for every integrally mixed
constructible Weil sheaf on \(Z\).  Every successive closed complement has
strictly smaller support complexity.
\end{theorem}

\begin{definition}[Normalization defect of a curve]
For a reduced finite-type curve \(C_0\) with finite normalization
\(\nu:\widetilde C_0\to C_0\), define
\[
 \delta(C_0)=\operatorname{length}_{\mathcal O_{C_0}}
      (\nu_*\mathcal O_{\widetilde C_0}/\mathcal O_{C_0}).
\]
The defect is finite, vanishes exactly on the normal locus, and is supported
on finitely many closed points.
\end{definition}

\begin{proposition}[Normalization reduction]
For a constructible Weil sheaf \(\mathcal F\) on a reduced curve, there is
a Frobenius-equivariant distinguished triangle
\[
 \mathcal F\longrightarrow R\nu_*L\nu^*\mathcal F
 \longrightarrow P\longrightarrow\mathcal F[1],
\]
where \(P\) has finite support.  Weight bounds for the normalized term and
the finite-support term imply the same bound for \(\mathcal F\).  Finite
and radicial transport preserve the triangle and all its weight predicates.
\end{proposition}

\begin{proposition}[Coefficient-filtration reduction]
Let
\(0=\mathcal F_0\subset\mathcal F_1\subset\cdots\subset
\mathcal F_m=\mathcal F\) be a finite Frobenius-stable filtration.  If the
desired \(R^nf_!\)-bound holds for every quotient
\(\mathcal F_a/\mathcal F_{a-1}\), it holds for \(\mathcal F\).  The same
statement holds for all-conjugates integral bounds.
\end{proposition}

\begin{theorem}[Target open--closed detection]
Let \(v:V_0\hookrightarrow Y_0\) be open with closed complement \(t:T_0\).
For a Cartesian base change of \(f_0:X_0\to Y_0\), compact-support base
change gives
\[
 v^*Rf_!K\xrightarrow{\sim}Rf_{V,!}v_X^*K,
 \qquad
 t^*Rf_!K\xrightarrow{\sim}Rf_{T,!}t_X^*K.
\]
If the two right-hand sides have integral upper weight \(w\), then so does
\(Rf_!K\).  The assertion iterates over every finite ordered locally closed
stratification of the target and retains the all-conjugates version.
\end{theorem}

\begin{theorem}[Composition transfer]
For separated finite-type maps
\(X_0\xrightarrow{h} Z_0\xrightarrow{g}Y_0\) and a constructible Weil sheaf
\(\mathcal F\), the Frobenius-equivariant spectral sequence
\[
 E_2^{a,b}=R^ag_!R^bh_!\mathcal F
          \Longrightarrow R^{a+b}(gh)_!\mathcal F
\]
is bounded.  If \(R^bh_!\mathcal F\) is integrally mixed of weights at most
\(w+b\), and its \(R^ag_!\)-image is integrally mixed of weights at most
\(w+a+b\), then
\[
 R^n(gh)_!\mathcal F
 \text{ is integrally mixed of weights at most }w+n.
\]
The same implication holds for all-conjugates integral mixedness.
\end{theorem}

\begin{theorem}[Generic projection dichotomy]
Let \(S_0\) be a smooth irreducible locally closed subscheme of
\(\mathbf A^1_{Y_0}\).  After replacing its scheme-theoretic image by a
dense open, one of the following alternatives holds.
\begin{enumerate}[label=\textup{\(\roman*\)}]
\item A dense open of \(S_0\) is separated and quasi-finite over the target.
\item A dense open of \(S_0\) is a smooth relative curve, and every pure
lisse coefficient of integral weight admits a generic controlled
compactification and trace retraction after finite radicial and etale base
change.
\end{enumerate}
The removed source and target complements have strictly smaller support
complexity.
\end{theorem}

\begin{theorem}[Relative affine-line projection]
For the projection
\(p:\mathbf A^1_{Y_0}\to Y_0\) and every
\(K\in\WeilD(\mathbf A^1_{Y_0},E)\),
\[
               \IntWtLe(K,w)\Longrightarrow
               \IntWtLe(Rp_!K,w).
\]
For an integrally mixed sheaf \(\mathcal F\) of weights at most \(w\),
\[
 R^rp_!\mathcal F
 \text{ is integrally mixed of weights at most }w+r.
\]
The quasi-finite branch uses the separated quasi-finite theorem, while the
one-dimensional branch uses the controlled relative-curve estimate.  The
all-conjugates statement holds when the affine-line and relative-curve
inputs have their all-conjugates forms.
\end{theorem}

\begin{proposition}[Closed-point detection for the projection]
There is a finite target stratification on which every
\(R^rp_!\mathcal F\) is lisse.  For a closed point \(y\) of a stratum,
compact-support base change gives an isomorphism of finite
\(Q_y\)-Frobenius modules
\[
 (Q_y,(R^rp_!\mathcal F)_{\bar y},F_y)
 \simeq
 (Q_y,H_c^r(\mathbf A^1_{\bar y},
              \mathcal F|_{\mathbf A^1_{\bar y}}),F_y),
 \qquad Q_y=q^{\deg(y)}.
\]
The finite filtrations constructed in the two branches restrict to the
stalk filtrations on the right.
\end{proposition}

\section{Affine targets and global assembly}

\begin{proposition}[Affine generator factorization]
Let \(V_0=\Spec B\) and let \(f_0:U_0=\Spec A\to V_0\) be of finite type.
A finite set of \(B\)-algebra generators gives a closed immersion
\[
                 U_0\xrightarrow{c}\mathbf A^N_{V_0}.
\]
The projection \(\mathbf A^N_{V_0}\to V_0\) is a composite of \(N\)
relative affine-line projections.  The factorization commutes with target
base change and scheme Frobenius, and \(c_!=c_*\) is exact and
weight-neutral.
\end{proposition}

\begin{theorem}[Affine morphism over an affine target]
For \(f_0:\Spec A\to\Spec B\) of finite type and
\(K\in\WeilD(\Spec A,E)\),
\[
                 \IntWtLe(K,w)\Longrightarrow
                 \IntWtLe(Rf_!K,w).
\]
The all-conjugates variant holds only when the relative affine-line theorem
is supplied with its all-conjugates local-Fourier and relative-curve inputs.
\end{theorem}

\begin{definition}[Ordered affine target partition]
Let \(V_0=\bigcup_{a=1}^mV_a\) be a finite affine open cover.  Define
locally closed affine strata
\[
 T_a=V_a\setminus\bigcup_{b<a}V_b,
\]
discarding empty terms and refining each difference by finitely many affine
locally closed subschemes.  The union of the first \(a\) strata is closed in
the union of the first \(a+1\).
\end{definition}

\begin{theorem}[Affine morphism over an arbitrary target]
Let \(f_0:U_0\to V_0\) be affine and of finite type.  Base change to every
member of an ordered affine target partition is an affine morphism between
affine schemes.  Compact-support base change and target open--closed
detection give
\[
                 \IntWtLe(K,w)\Longrightarrow
                 \IntWtLe(Rf_!K,w)
\]
for every \(K\in\WeilD(U_0,E)\).  The corresponding all-conjugates variant
has the same additional hypotheses as the relative affine-line theorem.
\end{theorem}

\begin{theorem}[Finite affine source refinement]
For every separated finite-type
\(f_0:X_0\to Y_0\), there are finite ordered locally closed partitions
\[
 Y_0=\bigsqcup_aV_a,
 \qquad X_0\times_{Y_0}V_a=\bigsqcup_bU_{ab},
\]
such that \(V_a\) and \(U_{ab}\) are affine and
\(U_{ab}\to V_a\) is affine of finite type.  All source localization
triangles and target base-change isomorphisms are Frobenius equivariant.
\end{theorem}

\begin{theorem}[Global sheaf theorem]
Let \(f_0:X_0\to Y_0\) be separated and of finite type.  If
\(\mathcal F\) is integrally mixed of weights at most \(w\), then
\[
 R^if_!\mathcal F
 \text{ is integrally mixed of weights at most }w+i
\]
for every \(i\).
\end{theorem}

\begin{theorem}[Deligne's upper-weight theorem, integral form]
Let \(f_0:X_0\to Y_0\) be separated and of finite type.  For every
\(K\in\WeilD(X_0,E)\) and \(w\in\mathbf R\),
\[
 \boxed{\qquad
 \IntWtLe(K,w)\Longrightarrow\IntWtLe(Rf_!K,w).
 \qquad}
\]
Equivalently, \(\mathcal H^i(Rf_!K)\) is integrally mixed of weights at
most \(w+i\) for every \(i\).
\end{theorem}

\begin{corollary}[Ordinary upper bound for integral input]
Under the hypotheses of the preceding theorem,
\[
        \IntWtLe(K,w)\Longrightarrow\WtLe(Rf_!K,w).
\]
No unconditional assertion with an arbitrary real-weight hypothesis
\(\WtLe(K,w)\) replaces the integral hypothesis in this note.
\end{corollary}

\section{Algebraically closed coefficients}

\begin{definition}[Finite-stage algebraically closed notation]
For the filtered pair category and Karoubi envelope defined in the
coefficient section, write
\[
 \WeilD(X_0,\Qlb)_{\mathrm{fin}}
 :=\WeilD(X_0,\Qlb)_{\mathrm{2-colim}}.
\]
Thus every object is represented by a triple \((E,K_E,e_E)\), where
\(E/\Ql\) is finite, \(K_E\in\WeilD(X_0,E)\), and \(e_E^2=e_E\).
\end{definition}

\begin{definition}[Weights at a finite-stage retract]
For a representative \((E,K_E,e_E)\), define its upper weight by applying
the finite-stage predicate to the image of every
\(\mathcal H^m(e_E)\) after one finite splitting extension.  Define lower
weight using the stagewise Verdier dual.  Integral and all-conjugates
integral predicates require the corresponding finite filtrations after one
common finite extension.
\end{definition}

\begin{theorem}[Common-stage descent]
Every finite family of objects, idempotents, morphisms, distinguished
triangles, truncations, Frobenius structures, and six-operation comparison
maps in \(\WeilD(-,\Qlb)_{\mathrm{fin}}\) is represented over one finite
extension of \(\Ql\).  The functors \(Lf^*,Rf_!,Rf_*,Rf^!\), tensor,
 internal Hom, and Verdier duality commute coherently with
further finite scalar extension on their displayed domains.
\end{theorem}

\begin{theorem}[Algebraically closed integral upper theorem]
The stagewise exceptional direct images descend to an exact functor
\[
 Rf_!: \WeilD(X_0,\Qlb)_{\mathrm{fin}}
       \longrightarrow\WeilD(Y_0,\Qlb)_{\mathrm{fin}}.
\]
For every object \(K\) and real \(w\),
\[
               \IntWtLe(K,w)\Longrightarrow\IntWtLe(Rf_!K,w).
\]
\end{theorem}

\section{Dual lower bounds and proper purity}

\begin{theorem}[Ordinary direct image, integral lower form]
For a separated finite-type morphism \(f_0:X_0\to Y_0\), strict Verdier
duality and the Frobenius-equivariant isomorphism
\(D_YRf_*K\simeq Rf_!D_XK\) give
\[
 \boxed{\qquad
 \IntWtGe(K,w)\Longrightarrow\IntWtGe(Rf_*K,w).
 \qquad}
\]
The theorem descends to the finite-stage \(\Qlb\)-category.
\end{theorem}

\begin{corollary}[Proper direct image]
If \(f_0\) is proper, the Frobenius-equivariant comparison
\(Rf_!\xrightarrow{\sim}Rf_*\) gives both implications
\[
 \IntWtLe(K,w)\Longrightarrow\IntWtLe(Rf_*K,w),
 \qquad
 \IntWtGe(K,w)\Longrightarrow\IntWtGe(Rf_*K,w).
\]
Thus an integrally pure complex of weight \(w\) has integrally pure proper
direct image of weight \(w\).
\end{corollary}

\begin{corollary}[Cohomology over a finite field]
Let \(a:X_0\to\Spec k\) be separated of finite type.  If
\(\IntWtLe(K,w)\), then the finite \(q\)-Frobenius module
\((q,H_c^i(X,K),F)\) is integrally
\(\MixLe_{q,\iota}\) of weights at most \(w+i\).  If
\(\IntWtGe(K,w)\), then
\((q,H^i(X,K),F)\) is integrally
\(\MixGe_{q,\iota}\) of weights at least \(w+i\).  If \(X_0\) is proper
and \(K\) is integrally pure of weight \(w\), then
\[
               \Pure_{q,\iota}(H^i(X,K),F;w+i)
\]
for every \(i\).
\end{corollary}

\begin{corollary}[Smooth proper families]
Let \(f_0:X_0\to Y_0\) be smooth and proper, assume smooth purity, and let
\(\mathcal L\) be lisse and
punctually pure of integral weight \(w\).  Then \(R^if_*\mathcal L\) is
lisse and punctually pure of weight \(w+i\).  At a closed point
\(y\in|Y_0|\), smooth proper base change identifies its stalk with the
finite \(Q_y\)-Frobenius module
\[
 (Q_y,H^i(X_{\bar y},\mathcal L|_{X_{\bar y}}),F_y),
 \qquad Q_y=q^{\deg(y)},
\]
which is \(\Pure_{Q_y,\iota}\) of weight \(w+i\).  The conclusion is
all-conjugates pure when the input is all-conjugates pure.
\end{corollary}

\section{References}

The principal sources are Deligne, \emph{La conjecture de Weil II},
Publ. Math. IHES 52 \(1980\), especially Theorem 3.3.1; Laumon,
\emph{Transformation de Fourier, constantes d'equations fonctionnelles et
conjecture de Weil}; and the Stacks Project chapters on etale cohomology,
pro-etale cohomology, compact support, and duality.

\end{document}
